\section{Trabalhos relacionados}

A literatura sobre inferência de modelos de \textit{machine learning} se
organiza em duas vertentes que este trabalho conecta: comparações entre
motores de inferência para \textit{ensembles} de árvores e comparações
entre plataformas de \textit{serving} sob carga concorrente. Cada vertente tem
contribuições maduras, mas tipicamente operadas em isolamento. A combinação
das duas, em particular sob protocolo de medição de modelo aberto, recebe
menos atenção empírica do que mereceria.

\subsection{Plataformas de inferência para ensembles de árvores}

Diversos trabalhos investigam o desempenho de motores especializados em
\textit{ensembles} de árvores de decisão. Hummingbird \cite{nakandala2020}
propõe traduzir o percurso condicional das árvores em operações de
multiplicação de matrizes (GEMM), permitindo execução em \textit{hardware}
de aceleração tensorial; o trabalho reporta uma faixa de latência da ordem
de centenas de microssegundos por linha em modelos com milhares de árvores,
com o sobrecusto de tensorização incluído na medida. Treelite
\cite{cho2018} compila o \textit{ensemble} para código C com dicas explícitas
de predição de salto, atingindo cerca de 28 microssegundos por linha em
\textit{hardware} Haswell. lleaves \cite{sieboehm2021} estende a ideia
empregando o compilador LLVM e atinge 9,6 microssegundos por linha no
mesmo \textit{benchmark}; no mesmo experimento, o ONNX Runtime nativo
foi medido em 11,0 microssegundos por linha, validando que a faixa de
poucas dezenas de microssegundos é alcançável sem compilação especializada.

Trabalhos voltados a recuperação de informação propõem estruturas de
travessia ainda mais agressivas. QuickScorer \cite{lucchese2015} usa
representações por vetor de bits para avaliar 1.000 árvores de profundidade
8 em 2,6 a 9 microssegundos por documento em processadores Xeon.
RapidScorer \cite{ye2018} estende a ideia maximizando compactação a fim
de explorar paralelismo SIMD, reportando latência da ordem de microssegundos
para tamanhos de \textit{ensemble} comparáveis. Por fim, Ye et al.
\cite{asadi2024} comparam Treelite, Hummingbird, TensorFlow Decision
Forests e ONNX Runtime sob a perspectiva de sistemas de banco de dados,
confirmando que ONNX Runtime é competitivo em lotes pequenos e que a
faixa de poucas dezenas de microssegundos por linha é o regime esperado para
modelos de tamanho moderado em CPU.

A limitação compartilhada por essa vertente é que todos esses trabalhos
medem o motor de inferência em isolamento, em laços apertados de
\textit{Python} com \textit{cache} quente e tensores pré-construídos. Eles
estabelecem o piso teórico da inferência, mas não capturam o sobrecusto da
camada de \textit{serving} ao redor dela: \textit{parsing} de requisições
HTTP, despacho de eventos do \textit{framework} web, formação de lotes em
nível de aplicação, agendamento de \textit{threads} e contenção de
\textit{thread pool}.

\subsection{Comparações entre plataformas de serving para ML}

Na vertente de plataformas de \textit{serving}, há comparações entre FastAPI,
Triton Inference Server, TensorFlow Serving e TorchServe. Zaharia et al.
\cite{zaharia2024} comparam FastAPI e Triton sobre Kubernetes em uma
aplicação de inferência médica e reportam diferenças em vazão e latência,
mas o protocolo adotado é de modelo fechado e a métrica de comparação se
apoia em médias e erros, ignorando a cauda. A documentação do BentoML
\cite{bentoml_dispatcher} descreve o desenho do \textit{dispatcher} de
\textit{adaptive batching}, mas não publica microbenchmarks isolados do
mecanismo de fila e despacho, deixando ao usuário a tarefa de calibrar
empiricamente os parâmetros sob a sua carga real.

A comparação entre Triton, TensorFlow Serving e TorchServe é predominante
em literatura voltada a redes profundas em GPU, onde o custo fixo por
chamada é elevado o suficiente para que o lote amortize-se. Este trabalho
deliberadamente fica fora desse regime, optando por um modelo tabular barato
em CPU para isolar o trade-off oposto: onde a infraestrutura de \textit{serving}
domina o tempo de resposta e onde a formação de lotes é mais cara do que a
inferência que ela tentaria amortizar.

\subsection{Posicionamento deste trabalho}

A contribuição empírica desta monografia preenche a lacuna entre as duas
vertentes. O motor de inferência usado (ONNX Runtime sobre o operador
\textit{TreeEnsembleClassifier}) é o mesmo medido por Sieboehm
\cite{sieboehm2021} e por Ye et al. \cite{asadi2024}, com piso de
inferência consistente com seus números. Sobre esse piso, a varredura de
modelo aberto via Vegeta \cite{vegeta2024} e a decomposição em três camadas
permitem atribuir cada milissegundo medido a um componente específico da
plataforma, e a coleta direta de utilização de CPU durante o ataque identifica
o gargalo arquitetural do \textit{dispatcher}. Nenhum dos trabalhos
revisados executa essa combinação de protocolo de modelo aberto, decomposição
de orçamento de latência e evidência de mecanismo em CPU para o regime
tabular barato.

